\documentclass[10pt,a4paper]{amsart}
\usepackage{amsmath,amssymb,amsthm,graphicx}
\usepackage{verbatim}
\newenvironment{reference}{\comment}{\endcomment}
%\newenvironment{reference}{}{}
\newenvironment{slogan}{\comment}{\endcomment}
\newenvironment{history}{\comment}{\endcomment}
% For commutative diagrams we use Xy-pic
\usepackage[all]{xy}

% We use 2cell for 2-commutative diagrams.
\xyoption{2cell}
% Theorem environments.
%
\theoremstyle{plain}
\newtheorem{theorem}[subsection]{Theorem}
\newtheorem{proposition}[subsection]{Proposition}
\newtheorem{lemma}[subsection]{Lemma}

\theoremstyle{definition}
\newtheorem{definition}[subsection]{Definition}
\newtheorem{example}[subsection]{Example}
\newtheorem{exercise}[subsection]{Exercise}
\newtheorem{situation}[subsection]{Situation}

\theoremstyle{remark}
\newtheorem{remark}[subsection]{Remark}
\newtheorem{remarks}[subsection]{Remarks}

\numberwithin{equation}{subsection}

% Macros
%
\def\lim{\mathop{\mathrm{lim}}\nolimits}
\def\colim{\mathop{\mathrm{colim}}\nolimits}
\def\Spec{\mathop{\mathrm{Spec}}}
\def\Hom{\mathop{\mathrm{Hom}}\nolimits}
\def\Ext{\mathop{\mathrm{Ext}}\nolimits}
\def\SheafHom{\mathop{\mathcal{H}\!\mathit{om}}\nolimits}
\def\SheafExt{\mathop{\mathcal{E}\!\mathit{xt}}\nolimits}
\def\Sch{\mathit{Sch}}
\def\Mor{\mathop{\mathrm{Mor}}\nolimits}
\def\Ob{\mathop{\mathrm{Ob}}\nolimits}
\def\Sh{\mathop{\mathit{Sh}}\nolimits}
\def\NL{\mathop{N\!L}\nolimits}
\def\CH{\mathop{\mathrm{CH}}\nolimits}
\def\proetale{{pro\text{-}\acute{e}tale}}
\def\etale{{\acute{e}tale}}
\def\QCoh{\mathit{QCoh}}
\def\Ker{\mathop{\mathrm{Ker}}}
\def\Im{\mathop{\mathrm{Im}}}
\def\Coker{\mathop{\mathrm{Coker}}}
\def\Coim{\mathop{\mathrm{Coim}}}

% Boxtimes
%
\DeclareMathSymbol{\boxtimes}{\mathbin}{AMSa}{"02}

%
% Macros for moduli stacks/spaces
%
\def\QCohstack{\mathcal{QC}\!\mathit{oh}}
\def\Cohstack{\mathcal{C}\!\mathit{oh}}
\def\Spacesstack{\mathcal{S}\!\mathit{paces}}
\def\Quotfunctor{\mathrm{Quot}}
\def\Hilbfunctor{\mathrm{Hilb}}
\def\Curvesstack{\mathcal{C}\!\mathit{urves}}
\def\Polarizedstack{\mathcal{P}\!\mathit{olarized}}
\def\Complexesstack{\mathcal{C}\!\mathit{omplexes}}
% \Pic is the operator that assigns to X its picard group, usage \Pic(X)
% \Picardstack_{X/B} denotes the Picard stack of X over B
% \Picardfunctor_{X/B} denotes the Picard functor of X over B
\def\Pic{\mathop{\mathrm{Pic}}\nolimits}
\def\Picardstack{\mathcal{P}\!\mathit{ic}}
\def\Picardfunctor{\mathrm{Pic}}
\def\Deformationcategory{\mathcal{D}\!\mathit{ef}}
\title{The stalk of the sheaf of homomorphisms computes the homomorphisms of the stalks for finitely presented modules}
\author{The Stacks Project Authors}
\date{stacks-project:modules.tex:local-snapshot-20261012; GNU FDL 1.2 or later}
\begin{document}
\maketitle
\noindent\textit{Editorial note.} All delivered mathematical text is the author's own, copied from the pinned Stacks Project chapter source sources/stacks-modules/pinned-original-source.tex; no statement or proof step was rewritten, reordered or removed.\par\medskip
\noindent\textit{Editorial note.} Preamble provenance: the Stacks Project keeps one single shared preamble for every chapter; this card takes it directly from the staged sources/stacks-modules/preamble.tex as three verbatim spans (lines 8-12, 14-18 and 154-217), so no derived file was written and nothing was added or redefined. The wrapper class and the base packages are supplied by the assembler.\par\medskip
\noindent\textit{Editorial note.} Bibliography provenance: neither the counted statement nor its proof contains a citation (this lane rescans the whole span before writing the card), so the card carries no bibliography and no entry was reconstructed.\par\medskip
\noindent\textit{Editorial note.} Selection provenance: this card was opened under the file-based control contract (CONTROL.md), with the mechanical pre-card gate run on its statement and proof interval before the card was written.\par\medskip
\noindent\textit{Allowed standard prerequisites (editorial).}\begin{itemize}
\item The exactness of the sheaf of homomorphisms in the first variable for a right exact sequence, and in the second variable for a left exact sequence. Author lemma lemma-internal-hom-exact, whose statement is reproduced verbatim as a context span of this card (source lines 3775-3799). Produces the exact sequence whose stalk is compared with the corresponding sequence of the stalks.
\item The description of the stalk of the sheaf of homomorphisms as classes of local homomorphisms, and local generation of a module of finite type by finitely many sections. Used inside the counted proof at source lines 3851-3863. Reduces injectivity to testing the given local homomorphism on finitely many generators.
\end{itemize}\medskip
\begin{lemma}
\label{lemma-stalk-internal-hom}
Let $(X, \mathcal{O}_X)$ be a ringed space.
Let $\mathcal{F}$, $\mathcal{G}$ be $\mathcal{O}_X$-modules.
If $\mathcal{F}$ is of finite type then the canonical map
$$
\SheafHom_{\mathcal{O}_X}(\mathcal{F}, \mathcal{G})_x
\to
\Hom_{\mathcal{O}_{X, x}}(\mathcal{F}_x, \mathcal{G}_x)
$$
is injective.
If $\mathcal{F}$ is finitely presented,
this canonical morphism is an isomorphism.
\end{lemma}

\begin{lemma}
\label{lemma-internal-hom-exact}
Let $(X, \mathcal{O}_X)$ be a ringed space.
Let $\mathcal{F}$, $\mathcal{G}$ be $\mathcal{O}_X$-modules.
\begin{enumerate}
\item If $\mathcal{F}_2 \to \mathcal{F}_1 \to \mathcal{F} \to 0$
is an exact sequence of $\mathcal{O}_X$-modules, then
$$
0 \to
\SheafHom_{\mathcal{O}_X}(\mathcal{F}, \mathcal{G}) \to
\SheafHom_{\mathcal{O}_X}(\mathcal{F}_1, \mathcal{G}) \to
\SheafHom_{\mathcal{O}_X}(\mathcal{F}_2, \mathcal{G})
$$
is exact.
\item If $0 \to \mathcal{G} \to \mathcal{G}_1 \to \mathcal{G}_2$
is an exact sequence of $\mathcal{O}_X$-modules, then
$$
0 \to
\SheafHom_{\mathcal{O}_X}(\mathcal{F}, \mathcal{G}) \to
\SheafHom_{\mathcal{O}_X}(\mathcal{F}, \mathcal{G}_1) \to
\SheafHom_{\mathcal{O}_X}(\mathcal{F}, \mathcal{G}_2)
$$
is exact.
\end{enumerate}
\end{lemma}

\begin{proof}
The map sends the equivalence class of
$(U, \varphi)$ in $\SheafHom_{\mathcal{O}_X}(\mathcal{F}, \mathcal{G})_x$,
where $x \in U \subset X$ is open and
$\varphi \in \Hom_{\mathcal{O}_U}(\mathcal{F}|_U, \mathcal{G}|_U)$,
to the induced map on stalks at $x$, namely
$\varphi_x : \mathcal{F}_x \to \mathcal{G}_x$.

\medskip\noindent
Suppose $\mathcal{F}$ is of finite type.
Pick a representative $(U, \varphi)$ of an element $\sigma$
in the kernel of the map, i.e., $\varphi_x = 0$.
Shrinking $U$ if necessary, choose sections
$s^1, \ldots, s^n \in \mathcal{F}(U)$ generating $\mathcal{F}|_U$.
Since $\varphi_x(s^i_x) = 0$
and we are dealing with a finite number of sections,
we can find an open neighborhood 
$V \subset U$ of $x$ such that $\varphi_V(s^i|_V)=0$
for all $i = 1, \ldots, n$. Since $s^i|_V$, $i = 1, \ldots, n$ generate
$\mathcal{F}|_V$ this means that $\varphi|_V = 0$.
Since $(U, \varphi)$ is equivalent to $(V, \varphi|_V)$
we conclude $\sigma = 0$ and injectivity of the map follows.

\medskip\noindent
Next, assume $\mathcal{F}$ is finitely presented.
By localizing on $X$ we may assume that $\mathcal{F}$ has a presentation
$$
\bigoplus\nolimits_{j = 1, \ldots, m}
\mathcal{O}_X
\longrightarrow
\bigoplus\nolimits_{i = 1, \ldots, n}
\mathcal{O}_X
\to
\mathcal{F}
\to
0.
$$
By Lemma \ref{lemma-internal-hom-exact} this gives an exact sequence
$
0 \to
\SheafHom_{\mathcal{O}_X}(\mathcal{F}, \mathcal{G}) \to
\bigoplus\nolimits_{i = 1, \ldots, n} \mathcal{G}
\longrightarrow
\bigoplus\nolimits_{j = 1, \ldots, m} \mathcal{G}.
$
Taking stalks we get an exact sequence
$
0 \to
\SheafHom_{\mathcal{O}_X}(\mathcal{F}, \mathcal{G})_x \to
\bigoplus\nolimits_{i = 1, \ldots, n} \mathcal{G}_x
\longrightarrow
\bigoplus\nolimits_{j = 1, \ldots, m} \mathcal{G}_x
$
and the result follows since $\mathcal{F}_x$ sits in
an exact sequence
$
\bigoplus\nolimits_{j = 1, \ldots, m}
\mathcal{O}_{X, x}
\longrightarrow
\bigoplus\nolimits_{i = 1, \ldots, n}
\mathcal{O}_{X, x}
\to
\mathcal{F}_x
\to
0
$
which induces the exact sequence
$
0 \to
\Hom_{\mathcal{O}_{X, x}}(\mathcal{F}_x, \mathcal{G}_x) \to
\bigoplus\nolimits_{i = 1, \ldots, n} \mathcal{G}_x
\longrightarrow
\bigoplus\nolimits_{j = 1, \ldots, m} \mathcal{G}_x
$
which is the same as the one above.
\end{proof}

\end{document}
